
En
effet, raisonnons par récurrence sur @@M0@@. Lorsque @@M1@@, c'est clair car @@M2@@ et @@M3@@ est @@M4@@-régulier pour tout @@M5@@. Lorsque @@M6@@, on a une suite exacte @@M7@@ où @@M8@@ est un @@M9@@-faisceau. Par 1.7 et l'hypothèse de récurrence on voit que si @@M10@@, alors @@M11@@ et @@M12@@ est @@M13@@-régulier. Donc, par 1.4 @@M14@@ est @@M15@@-régulier et @@M16@@ pour @@M17@@. Donc, @@M18@@, ce qui est @@M19@@ par 1.8 (i) car @@M20@@. Enfin, @@M21@@ est @@M22@@-régulier en vertu de 1.3 (i); d'où la première assertion. La seconde en résulte grâce à 1.8 (ii) et 1.10.

\textbf{1.12.} Soient @@M23@@ un schéma noethérien et @@M24@@ un @@M25@@-schéma de type fini. Soit @@M26@@ une \uline{famille de classes de faisceaux cohérents sur les fibres de} @@M27@@; c'est-à-dire, pour tout point @@M28@@ et toute extension @@M29@@ de @@M30@@, on se donne des faisceaux cohérents @@M31@@ sur le schéma algébrique @@M32@@ et on dit qu'un @@M33@@ et un @@M34@@ sont dans la même classe s'il existe des @@M35@@-homomorphismes de @@M36@@ dans une même extension @@M37@@ de @@M38@@ tels que @@M39@@ et @@M40@@ soient isomorphes sur @@M41@@. Alors la famille @@M42@@ est dite \uline{limitée par le faisceau cohérent} @@M43@@ \uline{sur} @@M44@@ où @@M45@@ est de type fini sur @@M46@@, si @@M47@@ est contenue dans la famille des classes des faisceaux @@M48@@ avec @@M49@@; on dit que @@M50@@ est \uline{limitée} s'il existe @@M51@@ comme dessus.

Supposons en plus que @@M52@@ soit projectif et muni d'un faisceau inversible (relativement) ample @@M53@@. Alors @@M54@@ est dite une @@M55@@-\uline{famille}, pour une suite d'entiers @@M56@@, si toute classe dans @@M57@@ est représentée par un @@M58@@ avec @@M59@@ algébriquement clos, qui est un @@M60@@-faisceau.

